\begin{onehalfspace}

Le sezioni precedenti studiano la QEC su circuiti dedicati e Shor sotto rumore fisico. Per collegare direttamente i risultati servirebbe un'intera esecuzione fault-tolerant: registri codificati, cicli di sindrome, decoder e operazioni non-Clifford protette. Questa architettura non viene simulata. Si studia invece la \textbf{sensibilit\`a} di Shor a un modello semplificato: dopo ogni gate si applica un errore Pauli indipendente, di probabilit\`a $p_L$, e si misura il successo della fattorizzazione. Il pedice ``L'' indica il ruolo concettuale del parametro. Non lo rende equivalente al tasso di errore logico misurato nei test dei codici.

% -----------------------------------------------
\section{L'architettura a quattro livelli}
\label{sec:quattro_livelli}

L'analisi separa Shor ideale, Shor con rumore fisico, QEC studiata in circuiti
isolati e proxy uniforme per gate. Questa distinzione evita di interpretare il
quarto livello come una simulazione fault-tolerant. Esistono infatti sia una
barriera dimensionale --- codificare i dodici qubit di $N=15$ conduce
rapidamente a centinaia di qubit --- sia una barriera strutturale: Stim tratta
efficientemente circuiti stabilizer, mentre l'aritmetica modulare e la
\ac{QFT} rendono Shor non-Clifford. Schema completo, strumenti e manifest sono
riportati nella Sezione~\ref{app:shor_quattro_livelli}.

% -----------------------------------------------
\section{Definizione operativa del proxy}
\label{sec:modello_logico}

In M8, $p_L$ \`e la \textbf{probabilit\`a totale di applicare un Pauli non
identit\`a dopo ciascun gate compilato incluso nel modello}. Per un gate su $q$
qubit il canale distribuisce $p_L$ uniformemente sui $4^q-1$ Pauli non
identit\`a. Poich\'e Qiskit Aer usa il parametro depolarizzante $\lambda_q$, lo
script imposta
\begin{equation}
    \lambda_1=\frac{4}{3}p_L,
    \qquad
    \lambda_2=\frac{16}{15}p_L.
\end{equation}
Il rumore agisce su \texttt{sx}/\texttt{x} e \texttt{cx}; le rotazioni
\texttt{rz} sono virtuali. $p_L$ \`e dunque un proxy \emph{per gate nella base
compilata}, non una probabilit\`a per ciclo QEC, round di sindrome o blocco
logico.

Il circuito canonico usa $N=15$, $a=7$, otto qubit di conteggio e quattro di
lavoro. Ogni punto aggrega 20 repliche indipendenti da 4096 shot; una misura ha
successo quando frazioni continue e controlli con il massimo comun divisore
restituiscono fattori non banali. Seed, conteggi dei gate e hash del circuito
sono conservati nella Sezione~\ref{app:shor_quattro_livelli}. La domanda
supportata \`e quindi come vari il successo di questa specifica implementazione
al crescere di un residuo Pauli uniforme, non quale distanza di surface code
corrisponda a ciascun punto.

% -----------------------------------------------
\section{Il contesto delle stime fault-tolerant}
\label{sec:resource_estimates}

Le stime per istanze crittografiche comprendono routing, cicli di sindrome,
fabbriche di stati magici e schedulazione e non derivano dall'estrapolazione di
$N=15$. La letteratura passa dai circa 20 milioni di qubit fisici stimati da
Gidney ed Eker\aa{} per RSA-2048~\cite{gidney2021} a proposte molto meno
costose, ma fondate su codici e assunzioni architetturali differenti. La
traiettoria completa, incluse le pi\`u recenti stime non ancora sottoposte a
revisione paritaria, \`e discussa nella Sezione~\ref{app:shor_risorse} e nella
Tabella~\ref{tab:traiettoria_risorse}. M8 resta pertanto una sensibilit\`a locale,
non una previsione di risorse hardware.

% -----------------------------------------------
\section{Risultati M8 e M13}
\label{sec:risultati_logico}

\subsection*{M8: successo per singola misura}

La Tabella~\ref{tab:m8_griglia_completa} riporta la griglia completa e gli
intervalli Wilson al 95\%. Il controllo di monotonia combina incertezza
binomiale e dispersione fra repliche con copertura simultanea sui confronti
adiacenti.

\begin{table}[H]
\centering
\small
\begin{tabular}{rcc@{\hspace{1.2cm}}rcc}
\toprule
$p_L$ & $P_{\mathrm{success}}$ & IC 95\% & $p_L$ & $P_{\mathrm{success}}$ & IC 95\% \\
\midrule
$0$ & 0,7489 & [0,7459;0,7518] & $0,005$ & 0,6450 & [0,6417;0,6483] \\
$0,0001$ & 0,7469 & [0,7439;0,7498] & $0,01$ & 0,5651 & [0,5617;0,5685] \\
$0,0005$ & 0,7373 & [0,7343;0,7403] & $0,02$ & 0,4518 & [0,4484;0,4552] \\
$0,001$ & 0,7250 & [0,7219;0,7280] & $0,05$ & 0,2953 & [0,2921;0,2984] \\
$0,0017$ & 0,7083 & [0,7051;0,7114] & $0,1$ & 0,2503 & [0,2474;0,2533] \\
$0,002$ & 0,6992 & [0,6961;0,7024] & $0,2$ & 0,2441 & [0,2412;0,2471] \\
 & & & $0,5$ & 0,2459 & [0,2430;0,2489] \\
\bottomrule
\end{tabular}
\caption{M8: probabilità di ricavare i fattori per singola misura, 20
repliche da 4096 shot per punto.}
\label{tab:m8_griglia_completa}
\end{table}

\begin{figure}[H]
\centering
\includegraphics[width=0.82\textwidth]{qec_shor_logico_curve.pdf}
\caption[Shor sotto un proxy fenomenologico per gate]{Probabilit\`a di successo
di Shor ($N=15$, $a=7$) in funzione di $p_L$, con IC Wilson al 95\%. Non sono
sovrapposti valori Steane o surface code, perch\'e le rispettive metriche
aggregate non sono direttamente convertibili nel proxy M8. Figura rigenerata
dal JSON schema 2 con \texttt{figure\_src/gen\_shor\_logico.py}.}
\label{fig:qec_shor_logico}
\end{figure}

Il controllo statistico non rileva aumenti significativi del successo al crescere del rumore, tenendo conto dei confronti multipli. La curva \`e quindi compatibile con una funzione non crescente, anche se non tutte le diminuzioni fra punti adiacenti sono significative. Il successo passa da 0,7489 a 0,4518 tra $p_L=0$ e $0,02$. Oltre $p_L=0,1$ si stabilizza vicino a 0,245. Questo livello dipende dagli esiti che la verifica classica di $N=15$ accetta come utili; non \`e un valore universale.

\subsection*{Ripetizioni indipendenti}

Il valore ideale 0,7489 riguarda una singola misura. Per tentativi indipendenti
la probabilit\`a cumulativa è
\begin{equation}
 P(k)=1-(1-P_s)^k.
 \label{eq:cumulativa}
\end{equation}

\begin{figure}[H]
\centering
\includegraphics[width=0.72\textwidth]{qec_shor_pk.pdf}
\caption[Probabilità cumulativa dai punti M8]{Probabilità cumulativa per
$p_L\in\{0;0,01;0,05;0,2\}$. La soglia dell'80\% richiede rispettivamente
2, 2, 5 e 6 tentativi; le bande propagano gli estremi degli intervalli Wilson.}
\label{fig:qec_shor_pk}
\end{figure}

Due tentativi danno circa 0,937 nel caso ideale e 0,811 per $p_L=0,01$.
L'80\% va quindi letto come soglia cumulativa, non come requisito del singolo
run; ipotesi e controllo statistico sono richiamati nella
Sezione~\ref{app:shor_m8_dettaglio}.

\subsection*{M13: TOP-4 con pareggi espliciti}

M13 usa lo stesso proxy in 200 repliche da 1024 shot per punto. Il TOP-4 resta
1,000 fino a $p_L=0,05$, bench\'e la massa utile per misura scenda da 0,7512 a
0,2979; a $p_L=0,2$ vale invece 0,670, con intervallo identificato
$[0,535;0,860]$ dovuto ai pareggi. Il contrasto preregistrato rispetto a
$p_L=0$ mostra una diminuzione di 0,330, IC Newcombe 95\%
$[0,266;0,398]$, superiore al margine 0,10. Protocollo, griglia e analisi di
sensibilit\`a ai pareggi sono conservati nella Sezione~\ref{sec:app_comp_topk}
dell'Appendice~\ref{app:compilazione}.

\subsection*{Barriera di scalabilit\`a e compatibilit\`a delle metriche}

M8 e M13 restano limitati a $N=15$. La variante Beauregard validata per $N=21$
raggiunge 21\,036 gate \texttt{cx} e profondit\`a 23\,081 su 22 qubit: il limite
osservato riguarda la simulazione classica rumorosa adottata, non l'algoritmo
quantistico. Inoltre errore Steane, errore logico del surface code e guadagno
del decoder hanno unit\`a statistiche diverse da $p_L$ e non possono essere
proiettati direttamente sulla curva. Il contratto completo \`e esplicitato
nella Sezione~\ref{app:shor_metriche} e nella
Tabella~\ref{tab:decoder_su_shor}.

% -----------------------------------------------
\subsection{Discussione}
\label{subsec:discussione_logico}

M8 stabilisce una relazione quantitativa entro il proprio modello: ridurre il
residuo Pauli per gate riporta il successo verso il limite ideale. Non dice
quale codice realizzi un determinato $p_L$, n\'e confronta Steane e surface
code a parit\`a di hardware. Servirebbero un gate set logico, il conteggio dei
cicli di correzione, latenze e correlazioni del decoder e operazioni
non-Clifford fault-tolerant.

Il confronto tra successo per misura e TOP-K chiarisce inoltre che quattro
candidati possono coprire rapidamente un esito utile nella piccola istanza
$N=15$, anche quando la massa utile \`e degradata. Un TOP-4 elevato non prova
quindi che l'informazione quantistica sia intatta e va letto insieme alla
metrica per misura e all'incertezza dei pareggi.

Esperimenti recenti su qubit logici confermano il verso qualitativo del
beneficio, ma non forniscono la conversione fra risultati QEC e M8. La loro
portata, la natura pre-compilata della dimostrazione di Shor e la distinzione da
altri esperimenti su atomi riconfigurabili sono documentate nella
Sezione~\ref{app:shor_metriche}. Perci\`o $N=15$ resta un'istanza didattica: la
fattibilit\`a crittografica \`e una domanda architetturale separata, dominata
dall'overhead fault-tolerant.

\section{Layout e readout: risultati delle verifiche M11 e M11b}
\label{sec:layout_risultati}

Le due verifiche usano i dati di calibrazione offline di \texttt{FakeSherbrooke} e si svolgono in simulazione. In M11 sono validi 40 dei 60 layout richiesti. Gli altri 20 hanno calibrazioni fuori dal dominio del modello e vengono esclusi, senza contarli come fallimenti di Shor. La correlazione fra punteggio aggregato di calibrazione e successo sui dati di verifica è $\rho_s=-0.184$, con IC bootstrap 95\% $[-0.456;0.124]$. Il campione non mostra quindi una capacità predittiva conclusiva del punteggio. Scegliere il layout con il punteggio minimo dà $P_{\mathrm{holdout}}=0.451$; scegliere quello con il massimo successo nei dati di selezione dà $0.516$. La differenza $+0.065$ è descrittiva e non ha un intervallo dedicato (Sezione~\ref{sec:app_comp_pilota}).

M11b considera nove dei dodici sottografi richiesti, con cinque strategie iniziali per ciascuno: in totale 45 osservazioni di verifica. Si misura l'errore di lettura dei qubit effettivamente misurati dopo il routing. Tenendo conto di sottografo, numero di ECR e profondità, la correlazione parziale con il successo è $\rho_s=-0.605$, con IC cluster bootstrap 95\% $[-0.826;-0.190]$. A un errore di lettura maggiore si associa quindi un successo minore. È un'associazione osservata in soli nove cluster: non dimostra che assegnare inizialmente i registri ai qubit con readout migliore causi un aumento del successo. Protocollo, esclusioni, figure e analisi completa sono nella Sezione~\ref{sec:app_comp_ruoli} dell'Appendice~\ref{app:fase2}.

Questi controlli rispondono alla domanda di prevedibilità: nel campione
una metrica aggregata non basta a ordinare il successo, mentre una
proprietà del circuito finale conserva un'associazione misurabile.
Nessuna delle due prove addestra un selettore ML di layout.

\end{onehalfspace}
